\pageoriginale

\clearpage
\SGAExposeBookmark{XVII}{Pinceles de Lefschetz: teorema de existencia}
\begin{center}
{\large SGA 7 --- EXPOSÉ XVII}\par
\vspace{2em}
{\Large\underline{Pinceles de Lefschetz: teorema de existencia}}\par
\vspace{1.5em}
{\large por N. M. Katz}
\end{center}

\section*{0. Introducción}

En la presente exposición (que, en el seminario oral, precedía a los
Exposés XIV y XV, consagrados a la fórmula de Picard--Lefschetz)
introducimos los ``pinceles de Lefschetz'' de secciones hiperplanas de
un esquema proyectivo y liso, y damos condiciones para su existencia.
En la exposición siguiente haremos de ellos un estudio cohomológico
elemental, basado en la estructura conocida de la cohomología de las
variedades explotadas. El hecho de que ciertas hipótesis sobre los
pinceles de Lefschetz que tendremos que introducir se satisfagan a
menudo resultará de la fórmula de Picard--Lefschetz del Exposé XV, de
naturaleza menos elemental que los resultados de la presente exposición y
del siguiente.

Ciertos desarrollos de la presente exposición y de la siguiente se expondrán
en otra parte, en un marco más general (esquemas de base generales,
hipótesis de cuasiproyectividad en lugar de proyectividad, etc.):
en particular, el estudio precohomológico de los pinceles de secciones
hiperplanas, y especialmente de los pinceles de Lefschetz, en EGA V,
y el estudio cohomológico de los esquemas explotados en

\pageoriginale

el Exposé VII de Jouanolou en SGA 5. El lector encontrará asimismo en
EGA V demostraciones en las que la persecución intrínseca de diagramas
reemplaza el uso de coordenadas, predilecto del redactor del presente
exposición.

Deseo agradecer a P. Deligne y P. Griffiths la ayuda que me prestaron
en la concepción de esta exposición. En toda ella, $k$ designa un cuerpo
algebraicamente cerrado, salvo mención expresa en contrario.

\section*{1. Recordatorio sobre las singularidades cuadráticas
(cf. VI, 6)}

\subsubsection*{1.1.}
Sea $Y$ un $k$-esquema puramente de dimensión
$\sisi{n}{n-1}$.\footnote{La impresión dice primero ``de dimensión
$n$'' y, en la presentación como hipersuperficie que sigue, impone
$n=1+\dim Y$. Esta última es la convención utilizada por la hessiana
y la forma normal mostradas.}
Un punto cerrado $y\in Y$ se llama una
\underline{singularidad cuadrática ordinaria} de $Y$ si:
\begin{enumerate}
\item[(a)] el completado $\widehat{\mathcal O}_{Y,y}$ del anillo local
de $Y$ en $y$ es isomorfo a
\[
k[[T_1,\ldots,T_n]]/(Q(T)),
\]
el cociente del anillo de series formales en
$n=1+\dim Y$ variables por el ideal engendrado por un elemento
$Q(T)$;
\item[(b)] $Q$ comienza en grado dos, y la subvariedad de
$\mathbb P^{n-1}_k$ definida por la anulación de la parte homogénea de
grado dos de $Q(T)$ es lisa sobre $k$ si $n\geq2$, mientras que la
parte de grado dos de $Q(T)$ es no nula si $n=1$.
\end{enumerate}

Si la característica de $k$ es distinta de dos, \underline{o} si $n$
es par, la condición (b) equivale a:

\pageoriginale

\begin{enumerate}
\item[(b$'$)] $Q$ comienza en grado dos y su matriz hessiana
\[
\left(
\frac{\partial^2Q}{\partial T_i\partial T_j}(0)
\right)_{1\leq i,j\leq n}
\]
es inversible.
\end{enumerate}
Esto se deduce del criterio jacobiano de lisitud.

La forma estándar de la parte homogénea de grado dos $Q_2$ de una
serie $Q$ que satisface (b), obtenible mediante un cambio lineal de
variables, es
\begin{equation}\tag{1.1.1}
\left\{
\begin{aligned}
Q_2(T)&=\sum_{i=1}^{n/2}T_{2i-1}T_{2i}
&&\text{si $n$ es par},\\[1ex]
Q_2(T)&=T_n^2+\sum_{i=1}^{(n-1)/2}T_{2i-1}T_{2i}
&&\text{si $n$ es impar}.
\end{aligned}
\right.
\end{equation}
Esto resulta de la teoría de las formas cuadráticas (Bourbaki,
\emph{Algèbre}, cap.~X).

\subsubsection*{1.2.}
Un punto cerrado $y\in Y$ que satisface (a) y (b) se llama una
\underline{singularidad cuadrática no degenerada}. Puesto que el
determinante de la matriz hessiana es idénticamente nulo en
característica dos cuando $n$ es impar, las singularidades cuadráticas
no degeneradas presentadas en $n$ variables solo existen si $n$ es
par o si la característica de $k$ es distinta de dos; en esos casos
son precisamente las singularidades cuadráticas ordinarias.

\pageoriginale

\section*{2. Pinceles de Lefschetz: enunciado de los resultados}

\subsubsection*{2.1.}
Denotemos por $\mathbb P^r$ el espacio proyectivo de dimensión $r$
sobre $k$, por $\check{\mathbb P}^{r}$ el espacio proyectivo
\underline{dual} de los hiperplanos de $\mathbb P^r$, y por
$\operatorname{Gr}(1,\check{\mathbb P}^{r})$ la variedad de rectas
de $\check{\mathbb P}^{r}$. Un punto
$D\in\operatorname{Gr}(1,\check{\mathbb P}^{r})$ se llama un
\underline{pincel de hiperplanos} de $\mathbb P^r$. Considerando $D$
como una recta de $\check{\mathbb P}^{r}$, escribimos, por abuso de
notación, $D=\{H_t\}_{t\in D}$, donde los $H_t$ son los hiperplanos
de $\mathbb P^r$ que constituyen los puntos de la recta $D$.

El \underline{eje de $D$} es la subvariedad lineal de codimensión dos
de $\mathbb P^r$ que es la intersección $H_0\cap H_\infty$ de dos
miembros distintos cualesquiera del pincel. Considerando este eje
como un punto de $\operatorname{Gr}(r-2,\mathbb P^r)$, la variedad de
subvariedades lineales de dimensión $r-2$ de $\mathbb P^r$, la
aplicación que envía $D$ a su eje es el conocido isomorfismo de
ortogonalidad
\[
\operatorname{Gr}(1,\check{\mathbb P}^{r})
\xrightarrow{\ \sim\ }
\operatorname{Gr}(r-2,\mathbb P^r).
\]

\subsubsection*{2.2.}
Sea $X$ un $k$-esquema propio, liso e irreducible, de dimensión
$n\geq1$, provisto de una inmersión
$X\hookrightarrow\mathbb P^r$. Un pincel
$D=\{H_t\}_{t\in\sisi{\mathbb P}{D}}$ de hiperplanos de
$\mathbb P^r$ se llama un \underline{pincel de Lefschetz} (con
respecto a esta inmersión) si:
\begin{enumerate}
\item[(a)] el eje de $D$ corta transversalmente a $X$
(EGA IV, 17.13.7);
\item[(b)] existe un abierto denso $U\subset D$ tal que, para todo
$t\in U$, el hiperplano $H_t$ corta transversalmente a $X$;
\item[(c)] para $t_0\in D-U$, el hiperplano $H_{t_0}$ corta
transversalmente a $X$ salvo en

\pageoriginale

un punto, que es un punto singular cuadrático ordinario de
$X\cap H_{t_0}$.
\end{enumerate}
Veremos más abajo, en 3.2.1, que los pinceles de Lefschetz forman un
\underline{abierto} de
$\operatorname{Gr}(1,\check{\mathbb P}^{r})$.

\subsubsection*{2.3.}
La inmersión $X\hookrightarrow\mathbb P^r$ se llama una
\underline{inmersión de Lefschetz} si los pinceles de Lefschetz
forman un abierto denso de
$\operatorname{Gr}(1,\check{\mathbb P}^{r})$.\footnote{Más
exactamente, forman un abierto denso en el conjunto de puntos cerrados
de $\operatorname{Gr}(1,\check{\mathbb P}^{r})$. En lo que sigue nos
permitiremos tales abusos de lenguaje.}

\subsubsection*{2.4.}
Recordemos ahora la noción de \underline{múltiplo de una inmersión
dada} $X\hookrightarrow\mathbb P^r$. Es el compuesto de esta inmersión
con la $d$-ésima inmersión
\sisi{de Segre}{de Veronese}\footnote{La impresión la llama inmersión
de Segre $d$-ésima; la aplicación monomial mostrada es la inmersión de
Veronese.}
\[
S_d:\mathbb P^r\hookrightarrow
\mathbb P^{\binom{r+d}{d}-1},
\qquad
(X_0,\ldots,X_r)\longmapsto(\ldots,X^W,\ldots),
\]
donde
\[
X^W=X_0^{W_0}\cdots X_r^{\sisi{W_0}{W_r}},
\qquad W_0+\cdots+W_r=d.
\]
Así, una ``sección de $X$ por una hipersuperficie de grado $d$'' para
la inmersión dada---en inglés, una \emph{hypersurface section of
degree $d$}---no es sino una

\pageoriginale

sección hiperplana para el $d$-ésimo múltiplo de dicha inmersión.

\subsubsection*{Teorema 2.5.}
\underline{Sea $X$ irreducible, propio y liso sobre $k$, de dimensión
$n\geq1$, provisto de una inmersión
$i:X\hookrightarrow\mathbb P^r$. Entonces:}
\begin{enumerate}
\item[2.5.1.] \underline{todo múltiplo $d$-ésimo de $i$, con
$d\geq2$, es una inmersión de Lefschetz;}
\item[2.5.2.] \underline{si $k$ tiene característica cero, la
inmersión dada $i$ es una inmersión de Lefschetz.}
\end{enumerate}
La demostración se divide en varias partes, que siguen.

\section*{3. La variedad dual}

\subsubsection*{3.1.}
Sea $I$ el ideal de $\mathbb P^r$ que define $X$. Denotemos por
$N=I/I^2$ el haz conormal, localmente libre de rango $r-n$ sobre $X$,
y por $\mathbb P(N)$ el fibrado proyectivo sobre $X$ formado por las
rectas de $\check N$. (Si $X=\mathbb P^r$, el fibrado
$\mathbb P(N)$ es vacío.) El fibrado $\mathbb P(N)$ se identifica con
la subvariedad de $X\times\check{\mathbb P}^{r}$ formada por los pares
$(x,H)$ tales que $H$ es tangente a $X$ en $x$, mediante la aplicación
\[
\nu:\mathbb P(N)\longrightarrow X\times\check{\mathbb P}^{r}.
\]
Más explícitamente,
\begin{equation}\tag{3.1.1}
\mathbb P(I/I^2)_x\ni F\longmapsto
\left(x,\text{ el hiperplano de ecuación }
\sum_i\frac{\partial F}{\partial X_i}(x)T_i=0\right),
\end{equation}
donde $F$ designa una sección local de $I/I^2$. Definimos
\begin{equation}\tag{3.1.2}
\varphi:\mathbb P(N)\longrightarrow\check{\mathbb P}^{r}
\end{equation}
como el compuesto
\[
\mathbb P(N)\xrightarrow{\nu}
X\times\check{\mathbb P}^{r}
\xrightarrow{\operatorname{pr}_2}\check{\mathbb P}^{r}.
\]

\pageoriginale

\subsubsection*{3.1.3.}
Denotemos por $\check X$ la imagen de $\varphi$. Es el conjunto de
hiperplanos que no cortan transversalmente a $X$, y se llama la
\underline{variedad dual} de $X$ para la inmersión
$i:X\hookrightarrow\mathbb P^r$.

\subsubsection*{Proposición 3.1.4.}
\underline{La variedad dual $\check X$ es irreducible y propia, de
dimensión a lo sumo $r-1$.}

\noindent\underline{Demostración.}
Es la imagen por $\varphi$ de $\mathbb P(N)$, que es irreducible de
dimensión $r-1$ si $X\ne\mathbb P^r$, y vacío si $X=\mathbb P^r$.

\subsubsection*{Observación 3.1.5.}
La interpretación de $\mathbb P(N)$ puede precisarse como sigue. Sea
$\mathcal Y\subset X\times\check{\mathbb P}^{r}$ el subesquema de
incidencia cuya fibra sobre todo punto
$H\in\check{\mathbb P}^{r}(k)$ es $X\cap H$. Es inmediato que
$\mathcal Y$ es liso de dimensión $n+r-1$, por ser un fibrado
proyectivo de dimensión relativa $r-1$ sobre $X$. Por tanto,
$\mathcal Y\to\check{\mathbb P}^{r}$ hace de $\mathcal Y$ una
intersección completa relativa de dimensión relativa virtual
$(n+r-1)-r=n-1$.

El \underline{subesquema jacobiano}
\[
J_{n-1}(\mathcal Y/\check{\mathbb P}^{r})
\]
(VI, 5.2) es precisamente $\mathbb P(N)$. La formación del subesquema
jacobiano de una intersección completa relativa es compatible con todo
cambio de base. En consecuencia, para toda recta
$D\subset\check{\mathbb P}^{r}$, el subesquema jacobiano
$J_{n-1}(\widetilde X/D)$, donde
\[
\widetilde X=\mathcal Y\times_{\check{\mathbb P}^{r}}D,
\]
es precisamente
$\mathbb P(N)\times_{\check{\mathbb P}^{r}}D$.

\subsubsection*{Proposición 3.2.}
\underline{El subconjunto $B\subset\check X$ formado por los
hiperplanos $H$ tales que $X\cap H$ tiene un único punto singular, y
que ese punto}

\pageoriginale

\underline{es una singularidad cuadrática ordinaria, es abierto.}

Este es un caso particular de XV, 1.3.2, basado en la teoría de los
``módulos henselianos'' de R. Elkik (XV, 1.1.4). El caso en que se
restringe a puntos cuadráticos \underline{no degenerados}---es decir,
cuando $\operatorname{car}(k)\ne2$ \underline{o}
$n=\dim X$ es par---también se deduce inmediatamente de 3.3, cuya
demostración no utiliza 3.2.

\subsubsection*{Corolario 3.2.1.}
\underline{Bajo las hipótesis de 3.2, sea $F=\check X-B$, un cerrado
de $\check X$. Entonces los pinceles de Lefschetz forman un abierto de
$\operatorname{Gr}(1,\check{\mathbb P}^{r})$, y las condiciones
siguientes son equivalentes:}
\begin{enumerate}
\item[3.2.2.] \underline{la inmersión dada
$i:X\hookrightarrow\mathbb P^r$ es una inmersión de Lefschetz;}
\item[3.2.3.] \underline{existe un pincel de Lefschetz de hiperplanos
para la inmersión dada;}
\item[3.2.4.] \underline{la codimensión de $F$ en
$\check{\mathbb P}^{r}$ es al menos dos.}
\end{enumerate}

\noindent\underline{Demostración.}
Una recta $D\subset\check{\mathbb P}^{r}$ es un pincel de Lefschetz si
y solo si:
\begin{enumerate}
\item[(3.2.5)] el eje de $D$ corta transversalmente a $X$;
\item[(3.2.6)] $D$ no está contenida en
$\sisi{X}{\check X}$;
\item[(3.2.7)] $D$ no corta a $F$.
\end{enumerate}

\pageoriginale

Las condiciones (3.2.5) y (3.2.6) definen abiertos de
$\operatorname{Gr}(1,\check{\mathbb P}^{r})$ y, como $F$ es cerrado,
lo mismo sucede con (3.2.7); por tanto, los pinceles de Lefschetz
forman un abierto. Puesto que $\dim X\leq r-1$, el abierto (3.2.6) es
no vacío y, por consiguiente, denso. El abierto (3.2.5) también es no
vacío y denso. En efecto, elijamos un hiperplano $H$ que corte
transversalmente a $X$; tal $H$ existe porque $\dim X\leq r-1$.
Como $\dim(X\cap H)\leq r-1$, existe un hiperplano $G$ que corta
transversalmente a $X\cap H$, y la recta $D$ que pasa por $H$ y $G$
tiene entonces un eje que corta transversalmente a $X$.

Así, la inmersión es de Lefschetz si y solo si el abierto (3.2.7) es
no vacío y, por tanto, denso. Como $F$ es cerrado, esto ocurre si y
solo si $F$ tiene codimensión al menos dos en
$\check{\mathbb P}^{r}$ (cf. EGA V, o Mumford,
\emph{Introduction to Algebraic Geometry}, p.~88, Cor.~3).

\subsubsection*{Corolario 3.2.8.}
\underline{Bajo las hipótesis de 3.2, supongamos que
$i:X\hookrightarrow\mathbb P^r$ es una inmersión de Lefschetz. Sea
$H_0$ un hiperplano que corta transversalmente a $X$, y sea
$\operatorname{Gr}(1,\check{\mathbb P}^{r})_{H_0}$ la subvariedad de
$\operatorname{Gr}(1,\check{\mathbb P}^{r})$ formada por las rectas
que pasan por $H_0$. Entonces los pinceles de Lefschetz que pasan por
$H_0$ forman un abierto denso de
$\operatorname{Gr}(1,\check{\mathbb P}^{r})_{H_0}$.}

\noindent\underline{Demostración.}
Por 3.2.1, dichos pinceles forman un abierto de
$\operatorname{Gr}(1,\check{\mathbb P}^{r})_{H_0}$. El argumento que
acabamos de dar para demostrar que los abiertos (3.2.5) y (3.2.6) son
no vacíos muestra que sus intersecciones con esta subvariedad también
son no vacías. Queda por demostrar que existen rectas que pasan por
$H_0$ y no cortan a $F$.

\pageoriginale

Para ello, sea $T\subset\check{\mathbb P}^{r}$ un hiperplano que no
contiene a $H_0$. Existe un isomorfismo
\begin{center}
\begin{tikzcd}[column sep=huge]
\operatorname{Gr}(1,\check{\mathbb P}^{r})_{H_0}
  & T \arrow[l,"\sim"']\\
\substack{\text{la recta que pasa por}\\ H_0\text{ y }\tau}
  & \tau \arrow[l,mapsto]
\end{tikzcd}
\end{center}
de modo que $\operatorname{Gr}(1,\check{\mathbb P}^{r})_{H_0}$
``es'' un espacio proyectivo de dimensión $r-1$. Las rectas de este
espacio que pasan por un punto de $F$ forman un cerrado de dimensión a
lo sumo $r-2$, pues son la imagen del morfismo
\begin{center}
\begin{tikzcd}[column sep=huge]
F \arrow[r] & \operatorname{Gr}(1,\check{\mathbb P}^{r})_{H_0}\\
f \arrow[r,mapsto] &
  \substack{\text{la recta que pasa por }H_0\\\text{y por }f}.
\end{tikzcd}
\end{center}

\subsubsection*{Proposición 3.3.}
\underline{El morfismo
$\varphi:\mathbb P(N)\to\check{\mathbb P}^{r}$ no está ramificado en
un punto cerrado $(x_0,H)$ si y solo si $x_0$ es una singularidad
cuadrática no degenerada (1.2) de $X\cap H$. En particular, el
subconjunto de $\mathbb P(N)$ formado por los puntos $(x,H)$ tales que
$x$ es una singularidad cuadrática no degenerada de $X\cap H$ es
abierto.}

\pageoriginale

\noindent\underline{Demostración de la Proposición 3.3.}
Podemos limitarnos al caso $X\ne\mathbb P^r$. Las dos propiedades cuya
equivalencia se afirma son independientes de la elección particular de
coordenadas, de modo que podemos describirlas mediante coordenadas
elegidas de una manera muy especial. Como $X$ es una subvariedad lisa
de $\mathbb P^r$ de dimensión $n\geq1$, elijamos coordenadas homogéneas
$X_0,\ldots,X_r$ tales que:
\begin{enumerate}
\item[3.3.1.] el punto $x_0\in\mathbb P^r$ sea
$(1,0,\ldots,0)$;
\item[3.3.2.] poniendo $x_i=X_i/X_0$ para $i=1,\ldots,r$, de modo que
$\widehat{\mathcal O}_{\mathbb P^r,x_0}$ se identifique con
$k[[x_1,\ldots,x_r]]$, el ideal que define la traza de $X$ en
$\operatorname{Spec}\widehat{\mathcal O}_{\mathbb P^r,x_0}$ esté
engendrado por elementos $g_i$ $(i=1,\ldots,r-n)$ que, en el
completado, tengan la forma
\[
g_i=x_{n+i}-h_i(x_1,\ldots,x_n),
\qquad i=1,\ldots,r-n,
\]
donde los $h_i$ son series únicamente en $x_1,\ldots,x_n$ y
pertenecen al cuadrado del ideal maximal de
$k[[x_1,\ldots,x_n]]$;
\item[3.3.3.] mediante (3.1.1), el hiperplano $H$ provenga de
$g_{\sisi{n-r}{r-n}}$.
\end{enumerate}

Pongamos $V=\operatorname{Spec}(\widehat{\mathcal O}_{X,x_0})$.
De 3.3.2 se deduce una trivialización
\[
(I/I^2)|_V\simeq\mathcal O_V^{r-n}
\]
dada por
\begin{equation}\tag{3.3.4}
(\Lambda_1,\ldots,\Lambda_{r-n})\longmapsto
\sum_{i=1}^{\sisi{n-r}{r-n}}\Lambda_i g_i.
\end{equation}
Junto con 3.3.3, esto proporciona un entorno $W$ de $(x_0,H)$

\pageoriginale

en $\mathbb P(N)|_V$ y un isomorfismo
$V\times\mathbb A^{r-n-1}\simeq W$, dado por
\begin{equation}\tag{3.3.5}
\begin{aligned}
V\times\mathbb A^{r-n-1}&\longrightarrow\mathbb P(I/I^2)|_V,\\
(v,(\lambda_1,\ldots,\lambda_{r-n-1}))&\longmapsto
\left(v,\left[g_{r-n}+
\sum_{i=1}^{r-n-1}\lambda_i g_i\right]\right),
\end{aligned}
\end{equation}
donde $\lambda_i=\Lambda_i/\Lambda_{r-n}$.

En particular, el completado del anillo local de $\mathbb P(N)$ en
$(x_0,H)$ se identifica con
\[
k[[x_1,\ldots,x_n,\lambda_1,\ldots,\lambda_{r-n-1}]].
\]
Expresemos ahora $\varphi$ en estas coordenadas, al nivel de los
completados:
\begin{equation}\tag{3.3.6}
\varphi(x,\lambda)
=\text{el hiperplano de ecuación }
\sum_{i=0}^{\sisi{n+r}{r}}\varphi_i(x,\lambda)T_i=0.
\end{equation}
Como el hiperplano imagen contiene al punto $x$, cuyas coordenadas
homogéneas son
\begin{equation}\tag{3.3.7}
x=(1,x_1,\ldots,x_n,h_1(x),\ldots,h_{r-n}(x)),
\end{equation}
se tiene
\begin{equation}\tag{3.3.8}
\varphi_0(x,\lambda)
=-\sum_{i=1}^{n}x_i\varphi_i(x,\lambda)
 -\sum_{j=1}^{r-n}h_j(x)\varphi_{n+j}(x,\lambda)
\end{equation}
en $\widehat{\mathcal O}_{\mathbb P(N),(x_0,H)}$. Además, por (3.1.1),
si
\[
D_i=\frac{\partial}{\partial x_i}\in
\operatorname{Der}_k(
\widehat{\mathcal O}_{\mathbb P^r,x_0},
\widehat{\mathcal O}_{\mathbb P^r,x_0}),
\]
entonces, para $1\leq i\leq r$,
\begin{equation}\tag{3.3.9}
\varphi_i(x,\lambda)
=D_i\!\left(g_{r-n}+
\sum_{j=1}^{r-n-1}\lambda_jg_j\right)(x)
\quad\text{en }\widehat{\mathcal O}_{\mathbb P(N),(x_0,H)}.
\end{equation}

\pageoriginale

Recordando que $g_i=x_{n+i}-h_i(x_1,\ldots,x_n)$, se obtiene
\begin{equation}\tag{3.3.10}
\varphi_i(x,\lambda)
=-\frac{\partial}{\partial x_i}
\left(h_{r-n}+\sum_{j=1}^{r-n-1}\lambda_jh_j\right),
\qquad 1\leq i\leq n,
\end{equation}
\begin{equation}\tag{3.3.11}
\sisi{\varphi_i(x,\lambda)=\lambda_i,\quad1\leq i\leq n}
{\varphi_{n+j}(x,\lambda)=\lambda_j,\quad1\leq j\leq r-n-1},
\end{equation}
y
\begin{equation}\tag{3.3.12}
\varphi_r(x,\lambda)=1.
\end{equation}
\footnote{La impresión da rangos incompatibles en (3.3.11). Las
coordenadas de (3.3.5), la fórmula de coeficientes (3.3.9) y el bloque
identidad de la matriz jacobiana impresa fuerzan los índices de la
rama corregida.}

Sea $\mathfrak m$ el ideal maximal de
$\widehat{\mathcal O}_{\mathbb P(N),(x_0,H)}$. Puesto que
$h_i\equiv0\pmod{\mathfrak m^2}$ para $i=1,\ldots,r-n$, se obtiene
\begin{equation}\tag{3.3.13}
\varphi_0(x,\lambda)\equiv0\pmod{\mathfrak m^2}
\end{equation}
y
\begin{equation}\tag{3.3.14}
\varphi_i(x,\lambda)
\equiv-\frac{\partial h_{r-n}}{\partial x_i}
\pmod{\mathfrak m^2},
\qquad i=1,\ldots,n.
\end{equation}
Reuniendo (3.3.11)--(3.3.14), la matriz jacobiana de $\varphi$ en
$(x_0,H)$ tiene la forma por bloques

\pageoriginale

\begin{equation}\tag{3.3.15}
J_{(x_0,H)}(\varphi)=
\left(
\begin{array}{c|c}
0_{1\times n} & 0_{1\times(r-n-1)}\\ \hline
-\operatorname{Hess}(h_{r-n})(0) & 0_{n\times(r-n-1)}\\ \hline
0_{(r-n-1)\times n} & I_{r-n-1}\\ \hline
0_{1\times n} & 0_{1\times(r-n-1)}
\end{array}
\right).
\end{equation}
Por tanto, $\varphi$ no está ramificado en $(x_0,H)$ si y solo si la
matriz hessiana
\[
\left(
\frac{\partial^2h_{r-n}}{\partial x_i\partial x_j}(0)
\right)_{1\leq i,j\leq n}
\]
es inversible. Por otra parte, 3.3.2 y 3.3.3 muestran que el completado
del anillo local de $X\cap H$ en $x_0$ es
\begin{equation}\tag{3.3.16}
k[[x_1,\ldots,x_n]]/
\bigl(h_{r-n}(x_1,\ldots,x_n)\bigr),
\end{equation}

\pageoriginale

y se concluye mediante la Definición 1.2.

\subsubsection*{Ejemplo 3.4.}
Supongamos que $X$ es una hipersuperficie de $\mathbb P^r$, definida
por una ecuación homogénea $F(X_0,\ldots,X_r)=0$. Entonces
$\mathbb P(N)=X$, y en coordenadas homogéneas
$\varphi:X\to\check{\mathbb P}^{r}$ viene dado por
\begin{equation}\tag{3.4.1}
\varphi(X_0,\ldots,X_r)=
\left(
\frac{\partial F}{\partial X_0},\ldots,
\frac{\partial F}{\partial X_r}
\right).
\end{equation}
Este es el \underline{morfismo de Gauss}.

He aquí un ejemplo en el que $\varphi$ está ramificado en todas partes.
Sea $p$ la característica de $k$, sea $r\geq1$, y tomemos la
hipersuperficie $X$ de ecuación
\begin{equation}\tag{3.4.2}
\sum_{i=0}^{r}X_i^{p^r+1}=0.
\end{equation}
Entonces
\begin{equation}\tag{3.4.3}
\varphi(X_0,\ldots,X_r)
=(X_0^{p^r},\ldots,X_r^{p^r}).
\end{equation}
Así, $\check X$ se reduce a un punto si $p=0$, mientras que si $p>0$
se tiene $\check X\simeq X$ y $\varphi$ es el morfismo de Frobenius.
De 3.3 (o de 3.5.0 más abajo) se deduce que, si $p>2$, la inmersión
considerada no es una inmersión de Lefschetz.

Volvamos al caso general.

\pageoriginale

\subsubsection*{Proposición 3.5.}
\underline{Supongamos que el morfismo
$\varphi:\mathbb P(N)\to\check{\mathbb P}^{r}$ es genéricamente no
ramificado (es decir, no está ramificado en todas partes). Entonces
$\varphi$ induce un morfismo birracional de $\mathbb P(N)$ sobre
$\check X$. Más precisamente, los tres abiertos siguientes de
$\sisi{X}{\check X}$ son idénticos:}
\begin{enumerate}
\item[(a)] \underline{el mayor abierto sobre el que $\check X$ es
liso;}
\item[(b)] \underline{el mayor abierto $V\subset\check X$ tal que
$\varphi$ induce un isomorfismo
$\varphi^{-1}(V)\xrightarrow{\sim}V$;}
\item[(c)] \underline{el abierto (cf. 3.2) formado por los
hiperplanos $H$ tales que $X\cap H$ tiene un único punto singular,
siendo este una singularidad cuadrática no degenerada.}
\end{enumerate}

\subsubsection*{Corolario 3.5.0.}
\underline{Supongamos que $\operatorname{car}(k)\ne2$ o que $\dim X$
es par. La inmersión dada es una inmersión de Lefschetz si y solo si
se cumple una de las dos condiciones siguientes:}
\begin{enumerate}
\item[(a)] $\varphi$ \underline{es genéricamente no ramificado;}
\item[(b)] $\dim\check X\leq r-2$.
\end{enumerate}

En efecto, la suficiencia resulta del criterio (3.2.4). En el caso
(a), la Proposición 3.5 identifica $F$ con el lugar singular de
$\check X$, por lo que
\[
\dim F<\dim\check X=r-1;
\]
en el caso (b), el criterio es inmediato. Recíprocamente, supongamos
que $D$ es un pincel de Lefschetz y que $\dim\check X=r-1$. Existe un
punto $H\in D\cap\check X$, y 3.3 junto con la Definición 2.2(c)
muestra que $\varphi$ no está ramificado en los puntos de
$\varphi^{-1}(H)$. Por tanto, es genéricamente no ramificado.

\noindent\underline{Demostración de la Proposición 3.5.}
La hipótesis implica
\[
\dim\check X=\dim\mathbb P(N)=r-1,
\]
y por 3.2 $\check X$ es una hipersuperficie irreducible de
$\check{\mathbb P}^{r}$.

\pageoriginale

Sea $R$ el lugar de ramificación de $\varphi$. Por hipótesis, es un
cerrado de $\mathbb P(N)$ de dimensión a lo sumo $r-2$; por tanto,
$\varphi(R)$ es un cerrado de $\check X$ de dimensión a lo sumo
$r-2$. Sea $V'$ el abierto denso
\begin{equation}\tag{3.5.1}
V'=(\check X-\varphi(R))\cap
   (\check X-\operatorname{Sing}(\check X))
\end{equation}
de $\check X$, y denotemos por $F$ su complementario:
\begin{equation}\tag{3.5.1.1}
F=\varphi(R)\cup\operatorname{Sing}(\check X).
\end{equation}

Demostremos primero que $\varphi$ induce un isomorfismo
$\varphi^{-1}(V')\simeq V'$ construyendo una inversa
$\psi:V'\to\varphi^{-1}(V')$. Consideremos
$\varphi^{-1}(V')$ inmerso en $X\times V'$ mediante (3.1.1), de modo
que $\varphi$ es la segunda proyección. Construir $\psi$ equivale
entonces a construir un morfismo $\psi_1:V'\to\mathbb P^r$ que haga
conmutativo el diagrama
\begin{equation}\tag{3.5.2}
\begin{tikzcd}[column sep=huge]
\varphi^{-1}(V') \arrow[r,"\varphi"] \arrow[d]
  & V' \arrow[d,"\psi_1"]\\
X \arrow[r,hook] & \mathbb P^r
\end{tikzcd}
\end{equation}
y poner después
$\psi(v)=(\psi_1(v),v)\in\mathbb P^r\times V'$.

Para toda hipersuperficie de un espacio proyectivo, el morfismo de
Gauss (3.4.1) está definido en cada punto liso. Apliquemos esto a la
hipersuperficie $\check X\subset\check{\mathbb P}^{r}$. Como $V'$
está contenido

\pageoriginale

en el lugar liso $V$ de $\check X$, el morfismo de Gauss proporciona
\begin{equation}\tag{3.5.3}
\psi_1:V'\longrightarrow
\check{\check{\mathbb P}}^{r}\simeq\mathbb P^r.
\end{equation}
Queda comprobar que (3.5.2) conmuta. Sea $(x_0,H_0)$ un punto cerrado
de $\varphi^{-1}(V')$. Debemos demostrar que $x_0$, considerado como
hiperplano de $\check{\mathbb P}^{r}$, es tangente a $\check X$ en
$H_0$; en otras palabras, que la tangencia de $H_0$ a $X$ en $x_0$
implica la tangencia de $x_0$ a $\check X$ en $H_0$.

Como $\varphi$ no está ramificado en $(x_0,H_0)$ y $\check X$ es liso
en $H_0=\varphi(x_0,H_0)$, su diferencial es un isomorfismo
\begin{equation}\tag{3.5.4}
d\varphi:
T_{\mathbb P(N),(x_0,H_0)}
\xrightarrow{\sim}T_{\check X,H_0}.
\end{equation}
La primera parte de la demostración queda entonces concluida por el
lema siguiente.

\subsubsection*{Lema 3.6 (Euler).}
\underline{Sea $(x_0,H_0)$ un punto cerrado de $\mathbb P(N)$, y sea
$T_L$ el fibrado tangente del hiperplano
$L\subset\check{\mathbb P}^{r}$ definido por $x_0$. Entonces, dentro
de $T_{\check{\mathbb P}^{r},H_0}$,}
\[
d\varphi\bigl(T_{\mathbb P(N),(x_0,H_0)}\bigr)
\subset T_{L,H_0}.
\]

\noindent\underline{Demostración.}
Elijamos coordenadas homogéneas $X_0,\ldots,X_r$ sobre $\mathbb P^r$
en las que $x_0$ sea el punto
$\underline A=(A_0,\ldots,A_r)$.

\pageoriginale

Un punto de $T_{X,x_0}$ es un punto de $\mathbb P^r$ con valores en
$k[\epsilon]/(\epsilon^2)$ de la forma
\begin{equation}\tag{3.6.1}
\underline A+\epsilon\underline B,
\end{equation}
tal que, para toda sección local $H$ de $I$,
\begin{equation}\tag{3.6.2}
\sum_{i=0}^{r}B_i
\frac{\partial H}{\partial X_i}(\underline A)=0.
\end{equation}
Elijamos además una sección local $H$ de $I$ que dé lugar a $H_0$
mediante (3.1.1). Como $\mathbb P(N)$ es localmente trivial sobre
$X$, un punto de $T_{\mathbb P(N),(x_0,H_0)}$ puede representarse por
\begin{equation}\tag{3.6.3}
(\underline A+\epsilon\underline B,\ H+\epsilon F),
\end{equation}
donde $F$ es una sección local de $I$ y $\underline A,\underline B$
son como arriba. Su imagen por $\varphi$ es el hiperplano sobre
$k[\epsilon]/(\epsilon^2)$ de ecuación
\begin{equation}\tag{3.6.4}
\sum_{i=0}^{r}
\frac{\partial(H+\epsilon F)}{\partial X_i}
(\underline A+\epsilon\underline B)T_i=0.
\end{equation}
Por tanto, la afirmación equivale a
\begin{equation}\tag{3.6.5}
\sum_{i=0}^{r}A_i
\frac{\partial(H+\epsilon F)}{\partial X_i}
(\underline A+\epsilon\underline B)=0.
\end{equation}
Como $\epsilon^2=0$, la fórmula de Taylor da
\begin{equation}\tag{3.6.6}
\frac{\partial(H+\epsilon F)}{\partial X_i}
(\underline A+\epsilon\underline B)
=\frac{\partial H}{\partial X_i}(\underline A)
+\epsilon\frac{\partial F}{\partial X_i}(\underline A)
+\epsilon\sum_{j=0}^{r}B_j
\frac{\partial^2H}{\partial X_j\partial X_i}(\underline A).
\end{equation}

\pageoriginale

Basta, pues, verificar las tres fórmulas siguientes:
\begin{equation}\tag{3.6.7}
\sum_{i=0}^{r}A_i
\frac{\partial H}{\partial X_i}(\underline A)=0,
\end{equation}
\begin{equation}\tag{3.6.8}
\sum_{i=0}^{r}A_i
\frac{\partial F}{\partial X_i}(\underline A)=0,
\end{equation}
y
\begin{equation}\tag{3.6.9}
\sum_{i,j=0}^{r}A_iB_j
\frac{\partial^2H}{\partial X_i\partial X_j}(\underline A)=0.
\end{equation}
Como $H$ y $F$ son homogéneos de grado cero, la fórmula de Euler da
\begin{equation}\tag{3.6.10}
\sum_iX_i\frac{\partial H}{\partial X_i}=0,
\qquad
\sum_iX_i\frac{\partial F}{\partial X_i}=0.
\end{equation}
Puesto que $\partial H/\partial X_j$ es homogéneo de grado $-1$, la
misma fórmula da
\[
\sum_iX_i\frac{\partial^2H}{\partial X_i\partial X_j}
=-\frac{\partial H}{\partial X_j}.
\]
En consecuencia, (3.6.9) se convierte en
\begin{equation}\tag{3.6.11}
-\sum_jB_j\frac{\partial H}{\partial X_j}(\underline A)=0,
\end{equation}
que es precisamente (3.6.2).

Esto demuestra que $\varphi$ induce un isomorfismo
\begin{equation}\tag{3.6.12}
\varphi^{-1}(V')\xrightarrow{\sim}V',
\end{equation}

\pageoriginale

cuya inversa es $\psi$. Observemos ahora que el morfismo
\begin{equation}\tag{3.6.13}
\left\{
\begin{aligned}
\psi &:V'\longrightarrow
  \mathbb P^r\times V'\subset\mathbb P^r\times\check X,\\
\psi(v)&=(\psi_1(v),v)
\end{aligned}
\right.
\end{equation}
se extiende a un morfismo
\begin{equation}\tag{3.6.14}
\left\{
\begin{aligned}
\widetilde\psi &:V\longrightarrow\mathbb P^r\times V,\\
\widetilde\psi(v)&=(\widetilde\psi_1(v),v),
\end{aligned}
\right.
\end{equation}
donde
\begin{equation}\tag{3.6.15}
\widetilde\psi_1:V\longrightarrow\mathbb P^r
\end{equation}
es el morfismo de Gauss, ahora considerado como definido sobre todo
$V$.

Acabamos de demostrar
\begin{equation}\tag{3.6.16}
\psi(V')\subset\mathbb P(N).
\end{equation}
Como $V'$ es denso en $V$ y $\mathbb P(N)$ es cerrado en
$\mathbb P^r\times\check X$, esto implica
\begin{equation}\tag{3.6.17}
\widetilde\psi(V)\subset\mathbb P(N).
\end{equation}
De hecho, es claro que incluso
\begin{equation}\tag{3.6.18}
\widetilde\psi(V)\subset\varphi^{-1}(V).
\end{equation}
De nuevo por continuidad, las identidades

\pageoriginale

\[
\psi\circ\varphi=\operatorname{id}_{\varphi^{-1}(V')},
\qquad
\varphi\circ\psi=\operatorname{id}_{V'}
\]
se extienden a
\[
\widetilde\psi\circ\varphi=\operatorname{id}_{\varphi^{-1}(V)},
\qquad
\varphi\circ\widetilde\psi=\operatorname{id}_{V}.
\]
Por tanto, los abiertos de 3.5(a) y 3.5(b) son idénticos.

Demostremos ahora que los abiertos de 3.5(a) y 3.5(c) son idénticos.
Por 3.3, este último es el subconjunto $U\subset\check X$ definido por
\begin{equation}\tag{3.6.19}
U=\left\{H\in\check X\ \middle|\
\begin{array}{l}
\varphi^{-1}(H)\text{ consta de un único punto, y}\\
\varphi\text{ no está ramificado en dicho punto}
\end{array}
\right\}.
\end{equation}
El esquema $\varphi^{-1}(U)$ es normal, pues es liso, y el morfismo
inducido
\[
\varphi:\varphi^{-1}(U)\longrightarrow U
\]
es birracional, como ya se ha demostrado; es radicial por la
definición de $U$, y propio. Es, por tanto, la normalización de $U$.
En particular, $U$ es geométricamente unirrama. Como este morfismo
no está ramificado, EGA IV, 18.10.1 muestra que es étale; EGA IV,
18.10.18 muestra entonces que es una inmersión abierta. En
consecuencia, es un isomorfismo y $U$ es liso. Esto demuestra
$3.5(c)\subset3.5(a)$.

\pageoriginale

La fórmula (3.6.19) hace evidente la inclusión
$3.5(b)\subset3.5(c)$, y la igualdad
$3.5(a)=3.5(b)$ concluye la demostración.

\subsubsection*{3.7.}
El Teorema 2.5 es verdadero cuando $\operatorname{car}(k)\ne2$, y
también cuando $X$ tiene dimensión par.

\noindent\underline{Demostración.}
Teniendo en cuenta 3.5.0 y 3.3, basta verificar la afirmación
siguiente:

\medskip
\noindent(3.7.1) Para todo punto cerrado $x_0\in X$ y todo entero
$d\geq2$, existe una hipersuperficie $H$ de grado $d$, tangente a
$X$ en $x_0$, tal que $x_0$ es una singularidad cuadrática ordinaria
de $X\cap H$.

Elijamos coordenadas homogéneas $X_0,\ldots,X_r$ tales que
$x_0=(1,0,\ldots,0)$ y que las funciones $X_i/X_0$,
$i=1,\ldots,n$, sean coordenadas locales sobre $X$ cerca de $x_0$.
Podemos tomar por $H$ la hipersuperficie de ecuación
\begin{equation}\tag{3.7.1.1}
H=X_0^{d-2}\sum_{i=1}^{n}X_i^2
\qquad\text{si }\operatorname{car}(k)\ne2,
\end{equation}
o
\begin{equation}\tag{3.7.1.2}
H=X_0^{d-2}\sum_{i=1}^{n/2}X_iX_{n/2+i}
\qquad\text{si }n\text{ es par}.
\end{equation}
Si $n$ es impar, bajo la primera hipótesis, la elección estándar
equivalente consignada en la fuente es
\begin{equation}\tag{3.7.1.3}
H=X_0^{d-2}
\left(X_n^2+
\sum_{i=1}^{\lfloor n/2\rfloor}
X_iX_{\lfloor n/2\rfloor+i}\right).
\end{equation}

\pageoriginale

Para la afirmación 2.5.2, observemos que o bien $\varphi$ es
genéricamente no ramificado, en cuyo caso se aplica 3.5.0, o bien está
ramificado en todas partes. En característica cero, este último caso
implica
$\dim\sisi{X}{\check X}\leq r-2$, de modo que el criterio (3.2.4) se
cumple automáticamente.\footnote{La impresión lleva
$\dim X\leq r-2$. La alternativa del Corolario 3.5.0 y el criterio
invocado requieren $\dim\check X\leq r-2$.}

\section*{4. El caso general}

Esbozaremos otro método de demostración de 2.5, válido en toda
característica, admitiendo el enunciado siguiente, próximo a 3.1 y
caso particular de XV, 1.1.4.

\medskip
\noindent(4.1) \underline{El subconjunto de $\mathbb P(N)$ formado
por los puntos $(x,H)$ tales que $x$ es una singularidad cuadrática
ordinaria de $X\cap H$ es abierto.}

Observemos que 4.1 se deduce de 3.3, salvo cuando
$\operatorname{car}(k)=2$ y $\dim X$ es impar.

\subsubsection*{4.1.1.}
Sea $R_1$ el cerrado complementario del abierto 4.1 en $\mathbb P(N)$,
y sea
\[
F_1=\varphi(R_1)\subset\check X\subset\check{\mathbb P}^{r}
\]
su imagen.

\subsubsection*{Lema 4.1.2.}
\underline{Si el abierto 4.1 es no vacío, entonces
$\operatorname{codim}_{\check{\mathbb P}^{r}}(F_1)\geq2$.}

\noindent\underline{Demostración.}
Por hipótesis,
\[
\dim R_1<\dim\mathbb P(N)=r-1,
\]
y $F_1=\varphi(R_1)$.

\subsubsection*{4.1.3.}
Sea $F_2$ la parte de $\check X$ formada por los hiperplanos tangentes
a $X$ en al menos dos puntos. Observemos que $F_2$ no es

\pageoriginale

cerrado, pero que
\begin{equation}\tag{4.1.4}
\overline{F_2}\subset F_1\cup F_2.
\end{equation}
Esto se debe a que $F_1\cup F_2=F$ (cf. 3.2) es cerrado. La inclusión
(4.1.4) ``dice'' que una \underline{confluencia} de dos singularidades
no es una singularidad cuadrática ordinaria. En todo caso, la
construcción de 4.2.4 mostrará que $F_2$ es constructible.

Para demostrar 2.5.1 aplicaremos el criterio 3.2.4. Necesitamos saber
que
\[
\operatorname{codim}_{\check{\mathbb P}^{r}}(F_1)\geq2
\quad\text{y}\quad
\operatorname{codim}_{\check{\mathbb P}^{r}}(F_2)\geq2.
\]
Por una parte, $\sisi{3.8.1}{3.7.1}$ y 4.1.2 muestran que, tras
reemplazar la inmersión dada por su múltiplo $d$-ésimo, $d\geq2$, se
tiene
\[
\operatorname{codim}_{\check{\mathbb P}^{r}}
(\sisi{F_2}{F_1})\geq2.
\]
\footnote{La impresión dice ``3.8.1'' y $F_2$ en esta frase. La
construcción cuadrática ordinaria es (3.7.1), y el Lema 4.1.2 se
aplica a $F_1$; la Proposición 4.2 proporciona la afirmación requerida
para $F_2$.}
Así, 2.5.1 quedará demostrado una vez conocida la proposición
siguiente.

\subsubsection*{Proposición 4.2.}
\underline{Sea $X$ una subvariedad lisa e irreducible de
$\mathbb P^r$, de dimensión $n\geq1$. Para todo entero $d\geq2$, la
parte $F_2$ de}
\[
\check{\mathbb P}^{N}
=\check{\mathbb P}^{\binom{\sisi{n}{r}+d}{d}-1}
\]
\underline{formada por las hipersuperficies de grado $d$ de
$\mathbb P^r$ tangentes a $X$ en al menos dos puntos tiene
codimensión al menos $2$ en $\check{\mathbb P}^{N}$.}
\footnote{La impresión lleva $\binom{n+d}{d}-1$ aquí y
$\binom{r+d}{d}-d$ en la Proposición 4.3. La inmersión de Veronese de
grado $d$ definida en 2.4 da
$N=\binom{r+d}{d}-1$ en ambos lugares.}

\subsubsection*{4.2.1.}
Comencemos por el caso especial $d=2$ y $X$ una subvariedad lineal
de $\mathbb P^r$. Basta encontrar una recta de
$\check{\mathbb P}^{N}$ que no corte a $F$ (cf. 3.2.4). Elijamos
coordenadas homogéneas $X_0,\ldots,X_r$ sobre $\mathbb P^r$ tales que
$X$ esté definida por

\pageoriginale

\[
X_{n+1}=\cdots=X_r=0,
\]
y coordenadas homogéneas $(\lambda,\mu)$ sobre $\mathbb P^1$.

Si $n$ es par, $n=2k$, tomamos el pincel de cuádricas
\begin{equation}\tag{4.2.2}
H_{(\lambda,\mu)}
=\lambda X_0^2+(\lambda X_k+\mu X_0)X_{2k}
 +\sum_{i=1}^{k-1}(\lambda X_i+\mu X_{i+1})X_{k+i}.
\end{equation}
Si $n$ es impar, $n=2k+1$, tomamos el pincel
\begin{equation}\tag{4.2.3}
H_{(\lambda,\mu)}
=(\lambda X_1+\mu X_0)X_0
 +\sum_{i=1}^{k}(\lambda X_{2i}+\mu X_i)X_{2i+1}.
\end{equation}
En ambos casos se comprueba que $X\cap H_{(\lambda,\mu)}$ es liso
para $(\lambda,\mu)\ne(0,1)$, mientras que
$X\cap H_{(0,1)}$ tiene un único punto singular, que además es una
singularidad cuadrática ordinaria. Por tanto, el pincel no corta a
$F$.

\subsubsection*{4.2.4.}
Sea $\operatorname{Btg}^{d}(X)$ la subvariedad de
\[
\check{\mathbb P}^{N}\times
\bigl((X\times X)-\Delta(X\times X)\bigr)
\]
formada por las ternas $(H,x_1,x_2)$ tales que $H$ es una
hipersuperficie de grado $d$ de $\mathbb P^r$, tangente a $X$ en
$x_1$ y en $x_2$, con $x_1\ne x_2$. La imagen de
$\operatorname{Btg}^{d}(X)$ por $\operatorname{pr}_1$ es el conjunto
$F_2$ de la Proposición 4.2, de modo que basta demostrar
\begin{equation}\tag{4.2.5}
\dim\operatorname{Btg}^{d}(X)\leq N-2
\quad\text{si $d\geq3$, o si $d=2$ y $X$ no es lineal.}
\end{equation}

Para ello, sea $\operatorname{Lin}(X\times X)$ la subvariedad de
$(X\times X)-\Delta(X\times X)$ dada por
\begin{equation}\tag{4.2.6}
\begin{aligned}
\operatorname{Lin}(X\times X)
&=\bigl\{(x_1,x_2)\in X\times X\mid x_1\ne x_2,\\[-1ex]
&\hspace{6em}\text{la recta engendrada por $x_1,x_2$ es tangente a
$X$ en $x_1$}\bigr\}\\
&=\bigl\{(x_1,x_2)\in X\times X\mid x_1\ne x_2,\\[-1ex]
&\hspace{6em}x_2\in X\cap\operatorname{Tg}_X(x_1)\bigr\}.
\end{aligned}
\end{equation}

\pageoriginale

Es evidente que:

\medskip
\noindent(4.2.7) \underline{Si $X$ no es una subvariedad lineal de
$\mathbb P^r$, entonces}
\[
\dim\operatorname{Lin}(X\times X)<2\dim X.
\]
Por consiguiente, para demostrar (4.2.5) basta mostrar que las fibras
del morfismo
\begin{equation}\tag{4.2.8}
\operatorname{Btg}^{d}(X)
\xrightarrow{\ \operatorname{pr}_{2,3}\ }
(X\times X)-\Delta(X\times X),
\end{equation}
consideradas como subvariedades de $\check{\mathbb P}^{N}$, tienen
\begin{equation}\tag{4.2.9}
\left\{
\begin{array}{ll}
\text{codimensión al menos $2\dim X+1$}
 &\text{sobre $\operatorname{Lin}(X\times X)$},\\[1ex]
\text{codimensión al menos $2\dim X+2$}
 &\text{sobre $(X\times X)-\Delta(X\times X)
                 -\operatorname{Lin}(X\times X)$.}
\end{array}
\right.
\end{equation}
La afirmación (4.2.9) se deduce de la proposición siguiente tomando
por $L_i$ el espacio tangente a $X$ en $x_i$, para $i=1,2$.

\subsubsection*{Proposición 4.3.}
\underline{Sean $x_1$ y $x_2$ dos puntos cerrados distintos de
$\mathbb P^r$, y sean $L_1$ y $L_2$ dos subespacios lineales de
$\mathbb P^r$ tales que $x_1\in L_1$ y $x_2\in L_2$. En el espacio}
\[
\check{\mathbb P}^{N}
=\check{\mathbb P}^{\binom{r+d}{d}-\sisi{d}{1}}
\]
\underline{de las hipersuperficies de grado $d$, el subespacio lineal
formado por las hipersuperficies tangentes a $L_1$ en $x_1$ y a
$L_2$ en $x_2$ tiene codimensión}
\[
\left\{
\begin{aligned}
\dim L_1+\dim L_2+2,
&\quad\text{si $d\geq3$, o si $d=2$ y $L_1\cap L_2$ no contiene}\\[-0.2ex]
&\quad\text{la recta $(x_1,x_2)$ que pasa por $x_1,x_2$},\\[1ex]
\dim L_1+\dim L_2+1,
&\quad\text{si $d=2$ y $L_1\cap L_2$ contiene la recta $(x_1,x_2)$.}
\end{aligned}
\right.
\]

\pageoriginale

\noindent\underline{Demostración.}
Para $d\geq3$ hay que demostrar que
$\dim L_1+\dim L_2+2$ condiciones lineales son independientes. Basta,
pues, tratar el caso $L_1=L_2=\mathbb P^r$, es decir, determinar
cuántas condiciones imponen a una hipersuperficie ser singular en
$x_1$ y en $x_2$. Elijamos coordenadas homogéneas tales que
\[
x_1=(1,0,\ldots,0),\qquad x_2=(0,1,0,\ldots,0).
\]
La hipersuperficie de ecuación $\sum_WA_WX^W=0$ es singular en
$(1,0,\ldots,0)$ si y solo si se anulan todos los coeficientes de los
monomios $X_0^{d-1}X_i$, $i=0,\ldots,r$; es singular en
$(0,1,0,\ldots,0)$ si y solo si se anulan todos los coeficientes de
los monomios $X_1^{d-1}X_i$, $i=0,\ldots,r$. Para $d\geq3$, todos
estos monomios son distintos, lo que demuestra la afirmación en este
caso.

Para $d=2$, el mismo cálculo da
\[
\operatorname{codim}\geq\dim L_1+\dim L_2+1,
\]
porque el monomio $X_0X_1$ se cuenta dos veces.

Supongamos primero que $L_1\cap L_2$ no contiene la recta
$(x_1,x_2)$. Nos reducimos al caso $L_1=\mathbb P^r$ y $L_2$ un
hiperplano de $\mathbb P^r$ que no contiene $(x_1,x_2)$. Elijamos
coordenadas homogéneas $X_0,\ldots,X_r$ tales que
\[
x_1=(1,0,\ldots,0),\qquad
\sisi{x_1}{x_2}=(0,1,0,\ldots,0),\qquad L_2=(X_0=0).
\]
Una cuádrica $\sum_{i\leq j}A_{ij}X_iX_j$ es singular en
$(1,0,\ldots,0)$ si y solo si $A_{0j}=0$ para $j=0,\ldots,r$, y es
tangente a

\pageoriginale

$X_0=0$ en $(0,1,\ldots,0)$ si y solo si
$A_{\sisi{i}{1}j}=0$ para $j=1,\ldots,r$. En conjunto son
\[
2r+1=\dim L_1+\dim L_2+2
\]
condiciones independientes.

Finalmente, supongamos que $L_1\cap L_2$ contiene la recta
$(x_1,x_2)$. Nos reducimos al caso $L_1=L_2=(x_1,x_2)$. En
coordenadas homogéneas con
\[
x_1=(1,0,\ldots,0),\qquad x_2=(0,1,0,\ldots,0),
\]
la cuádrica $\sum_{i\leq j}A_{ij}X_iX_j$ pasa por
$(1,0,\ldots,0)$ y es tangente allí a $L_1$ si y solo si
$A_{00}=A_{01}=0$. Pasa por $(0,1,0,\ldots,0)$ y es tangente allí a
$L_2$ si y solo si $A_{11}=A_{01}=0$. Por tanto, solo hay
\[
3=\dim L_1+\dim L_2+1
\]
condiciones.

\section*{5. El grado de la variedad dual por un método
``elemental''}

(Cf. XVIII, 3.2 para otro método cuando $\dim X$ es par
\underline{o} $\operatorname{car}(k)\ne2$.)

Seguimos utilizando las notaciones de 2.2; por tanto,
$X\hookrightarrow\mathbb P^r$ es una subvariedad lisa e irreducible
de dimensión $n$.

\subsubsection*{5.1.}
Consideremos primero el morfismo (3.1.1)
\begin{equation}\tag{5.1.0}
\nu:\mathbb P(N)\longrightarrow X\times\check{\mathbb P}^{r}
\end{equation}
desde un punto de vista más funtorial. Por definición, el haz conormal
$\check N$ figura en la sucesión exacta

\pageoriginale

de haces coherentes sobre $X$
\begin{equation}\tag{5.1.1}
0\longrightarrow\check N
\longrightarrow\Omega^1_{\mathbb P^r/k}|_X
\longrightarrow\Omega^1_X
\longrightarrow0.
\end{equation}
Por otra parte, escribiendo las formas diferenciales de $\mathbb P^r$
en coordenadas homogéneas, se obtiene una sucesión exacta de haces
sobre $\mathbb P^r$
\begin{equation}\tag{5.1.2}
0\longrightarrow\Omega^1_{\mathbb P^r/k}
\longrightarrow\mathcal O_{\mathbb P^r}(-1)^{r+1}
\longrightarrow\mathcal O_{\mathbb P^r}
\longrightarrow0.
\end{equation}
Su restricción a $X$ es, por tanto, la sucesión exacta
\begin{equation}\tag{5.1.3}
0\longrightarrow\Omega^1_{\mathbb P^r/k}|_X
\longrightarrow\mathcal O_X(-1)^{r+1}
\longrightarrow\mathcal O_{\sisi{\mathbb P^r}{X}}
\longrightarrow0.
\end{equation}
\footnote{La impresión conserva el último término como
$\mathcal O_{\mathbb P^r}$ después de restringir la sucesión a $X$;
el término restringido es $\mathcal O_X$.}
Combinando (5.1.1) y (5.1.3) se obtiene una inclusión
\begin{equation}\tag{5.1.4}
\check N\lhook\joinrel\longrightarrow\mathcal O_X(-1)^{r+1},
\end{equation}
que, por transposición, da una sobreyección
\begin{equation}\tag{5.1.5}
\mathcal O_X(1)^{r+1}\longrightarrow N.
\end{equation}
Se obtiene así una inmersión cerrada
\begin{equation}\tag{5.1.6}
\alpha:\mathbb P(N)
\longrightarrow\mathbb P\bigl(\mathcal O_X(1)^{r+1}\bigr)
\simeq X\times\check{\mathbb P}^{r},
\end{equation}
que es precisamente (5.1.0).

\pageoriginale

\subsubsection*{5.2. La idea del cálculo.}

Supongamos desde ahora que
$\check X=\varphi(\mathbb P(N))$ tiene dimensión $r-1$, y es por tanto
una hipersuperficie de $\check{\mathbb P}^{r}$. El cálculo de su grado
se hace naturalmente introduciendo las clases de Chern [1].

Para toda variedad proyectiva lisa $T$, denotemos por
$\operatorname{CH}(T)$ su anillo de Chow, graduado por la codimensión.
Para $z\in\operatorname{CH}(T)$, escribimos
\[
\int_Tz
\]
para el grado de su componente de dimensión cero. La clase en
$\operatorname{CH}^1(T)$ de un haz inversible $\mathcal L$ sobre $T$
se denotará también por $\mathcal L$ cuando no pueda surgir confusión
(por ejemplo, bajo el signo de integral).

Sea $L$ la clase de una recta en
$\operatorname{CH}(\check{\mathbb P}^{r})$. Entonces
\begin{equation}\tag{5.2.1}
\deg(\check X)
=\int_{\check{\mathbb P}^{r}}L\cdot\check X.
\end{equation}
Conjuntísticamente, $\check X$ es la imagen de $\mathbb P(N)$ por
$\varphi$ (3.1.2). Sin embargo, desde el punto de vista de la imagen
directa propia
\[
\varphi_*:\operatorname{CH}(\mathbb P(N))
\longrightarrow\operatorname{CH}(\check{\mathbb P}^{r}),
\]
se tiene
\[
\varphi_*[\mathbb P(N)]=\deg(\varphi)[\check X].
\]

\pageoriginale

Por consiguiente, (5.2.1) se convierte en
\begin{equation}\tag{5.2.2}
\deg(\check X)
=\frac{1}{\deg(\varphi)}
\int_{\check{\mathbb P}^{r}}L\cdot\varphi_*[\mathbb P(N)].
\end{equation}
Aquí, por abuso de notación, $\varphi$ en $\deg(\varphi)$ designa el
morfismo $\mathbb P(N)\to\check X$ inducido por (3.1.2); su grado es un
entero no negativo, tomado igual a cero precisamente cuando
$\dim\sisi{X}{\check X}<r-1=\dim\mathbb P(N)$.\footnote{La impresión
lleva $\dim X<r-1$. La condición para que el morfismo inducido hacia la
variedad dual sea genéricamente finito es $\dim\check X=r-1$.}
La fórmula de proyección reescribe (5.2.2) como
\begin{equation}\tag{5.2.3}
\deg(\check X)
=\frac{1}{\deg(\varphi)}
\int_{\mathbb P(N)}\varphi^*(L).
\end{equation}

Sea $\eta\in\operatorname{CH}(\check{\mathbb P}^{r})$ la clase de un
hiperplano. Como $L=\eta^{r-1}$, la ecuación (5.2.3) se convierte en
\begin{equation}\tag{5.2.4}
\deg(\check X)
=\frac{1}{\deg(\varphi)}
\int_{\mathbb P(N)}\bigl(\varphi^*(\eta)\bigr)^{r-1}.
\end{equation}
Calculemos ahora concretamente $\varphi^*(\eta)$.

Sean $\xi$ (respectivamente, $\tau$) el haz inversible canónico
$\mathcal O_{\mathbb P(N)}(1)$ sobre $\mathbb P(N)$ (respectivamente,
$\mathcal O_{\mathbb P(\mathcal O_X(1)^{r+1})}(1)$ sobre
$\mathbb P(\mathcal O_X(1)^{r+1})$). Por funtorialidad, (5.1.6) da
\begin{equation}\tag{5.2.5}
\alpha^*(\tau)=\xi.
\end{equation}

\pageoriginale

Bajo la identificación
\[
\mathbb P\bigl(\mathcal O_X(1)^{r+1}\bigr)
\simeq X\times\check{\mathbb P}^{r},
\]
y escribiendo $H=\mathcal O_X(1)$, se tiene evidentemente
\begin{equation}\tag{5.2.6}
\tau=\operatorname{pr}_1^*(H)\otimes\operatorname{pr}_2^*(\eta).
\end{equation}
Sea $\pi:\mathbb P(N)\to X$ el morfismo estructural. Entonces
\[
\pi=\operatorname{pr}_1\circ\alpha,\qquad
\varphi=\operatorname{pr}_2\circ\alpha.
\]
Aplicando $\alpha^*$ a (5.2.6) y utilizando (5.2.5), se obtiene
\begin{equation}\tag{\sisi{5.2.6}{5.2.6$'$}}
\xi=\alpha^*(\tau)=\pi^*(H)\otimes\varphi^*(\eta).
\end{equation}
\footnote{La impresión numera tanto esta fórmula como la precedente
(5.2.6); la prima distingue la segunda aparición en la rama
corregida.}
Así, en $\operatorname{CH}(\mathbb P(N))$,
\begin{equation}\tag{5.2.7}
\varphi^*(\eta)=\xi-\pi^*(H),
\end{equation}
y (5.2.3) se convierte en
\begin{equation}\tag{5.2.8}
\deg(\check X)
=\frac{1}{\deg(\varphi)}
\int_{\mathbb P(N)}\bigl(\xi-\pi^*(H)\bigr)^{r-1}.
\end{equation}
El segundo miembro de (5.2.8) se calcula mediante el lema siguiente
(cf. 5.3.3).

\subsubsection*{Lema 5.3.}
\underline{Sean $X$ una variedad cuasiproyectiva lisa y conexa de
dimensión $n$, y $M$ un haz localmente libre de rango finito $a$
sobre $X$. Denotemos por $\xi$ el haz inversible
$\mathcal O_{\mathbb P(M)}(1)$ sobre $\mathbb P(M)$,}

\pageoriginale

\underline{y por $\pi:\mathbb P(M)\to X$ el morfismo estructural.
Recordemos que $\operatorname{CH}(\mathbb P(M))$ es un
$\operatorname{CH}(X)$-módulo libre, vía $\pi^*$, con base
$1,\xi,\ldots,\xi^{a-1}$, donde $\xi$ satisface}
{\small
\begin{equation}\tag{5.3.1}
\xi^a-C_1\xi^{a-1}+C_2\xi^{a-2}-C_3\xi^{a-3}+\cdots=0,
\end{equation}
}
\underline{siendo los $C_i$ las clases de Chern de $M$ en
$\operatorname{CH}(X)$. Pongamos
$C(M)=1+C_1+C_2+\cdots$, la clase de Chern total de $M$; es un
elemento inversible de $\operatorname{CH}(X)$. Entonces:}
\begin{enumerate}
\item[(i)] \underline{se tiene}
\begin{equation}\tag{5.3.2}
\pi_*\!\left(\frac{\xi^{a-1}}{1+\xi}\right)
=\frac{1}{C(M)};
\end{equation}
\item[(ii)] \underline{sea
$F(T)\in\operatorname{CH}(X)[T]$ un polinomio tal que el coeficiente
de $T^i$ sea homogéneo de grado $n+a-1-i$. Entonces $F(\xi)$ tiene
grado máximo $n+a-1=\dim\mathbb P(M)$ en
$\operatorname{CH}(\mathbb P(M))$. Se tiene}
\begin{equation}\tag{5.3.3}
\int_{\mathbb P(M)}F(\xi)
=(-1)^{a-1}\int_X\frac{F(-1)}{C(M)}.
\end{equation}
\end{enumerate}

\noindent\underline{Demostración.}
Comencemos deduciendo (5.3.3) de (5.3.2). Escribamos
\begin{equation}\tag{5.3.4}
\frac{1}{C(M)}=\sum_{j\geq0}b_j,\qquad
b_0=1,\qquad b_j\in\operatorname{CH}^j(X),\qquad
b_j=0\quad\text{para }j<0.
\end{equation}

\pageoriginale

Entonces (5.3.2) equivale a
\begin{equation}\tag{5.3.5}
\pi_*(\xi^{a-1+j})=(-1)^jb_j
\qquad\text{siempre que }a-1+j\geq0,
\end{equation}
pues $\pi_*(\xi^q)=0$ para $q<a-1$ por razones de dimensión.

Por linealidad, basta demostrar (5.3.3) cuando
$F(T)=uT^{a-1+j}$, con $u\in\operatorname{CH}^{n-j}(X)$ homogéneo.
Por (5.3.5) y la fórmula de proyección,
\begin{equation}\tag{5.3.6}
\begin{aligned}
\int_{\mathbb P(M)}F(\xi)
&=\int_{\mathbb P(M)}\pi^*(u)\xi^{a-1+j}\\
&=\int_Xu\,\pi_*(\xi^{a-1+j})
=(-1)^j\int_Xu\,b_j.
\end{aligned}
\end{equation}
Por otra parte,
\begin{equation}\tag{5.3.7}
\begin{aligned}
(-1)^{a-1}\int_X\frac{F(-1)}{C(M)}
&=(-1)^{a-1}\int_X
  \frac{(-1)^{a-1+j}u}{C(M)}\\
&=(-1)^j\int_X\sum_i u\,b_i.
\end{aligned}
\end{equation}
Como $\deg u=n-j$ y $\deg b_i=i$, se tiene
\[
\int_Xu\,b_i=0\qquad\text{para }i\ne j.
\]
Así, (5.3.7) se reduce a
\begin{equation}\tag{5.3.8}
(-1)^{a-1}\int_X\frac{F(-1)}{C(M)}
=(-1)^j\int_Xu\,b_j,
\end{equation}
lo que demuestra la implicación anunciada.

Queda demostrar (5.3.2), o equivalentemente (5.3.5).

\pageoriginale

Por razones de dimensión,
\[
\pi_*(\xi^j)=0\qquad\text{si }j<a-1,
\]
mientras que
\[
\pi_*(\xi^{a-1})=1_X
\]
por la razón geométrica de que $\xi^{a-1}$ corta cada fibra en un
punto. La fórmula de proyección da, por tanto,
\begin{equation}\tag{5.3.9}
\sisi{\displaystyle
\pi_*\!\left(\sum_{j=0}^{a-1}\pi^*(U_i)\xi^i\right)=U_{j-i}}
{\displaystyle
\pi_*\!\left(\sum_{i=0}^{a-1}\pi^*(U_i)\xi^i\right)=U_{a-1}}.
\end{equation}
\footnote{La fórmula impresa mezcla los índices $i,j$ y lleva
$U_{j-i}$ en el segundo miembro. Las dos identidades de imagen directa
inmediatamente precedentes fuerzan el coeficiente $U_{a-1}$.}
Por (5.3.9), la afirmación (5.3.5) resulta del lema siguiente, en el
que $\xi$ se reemplaza por $T$ y $C_i$ por $S_i$, teniendo en cuenta
(5.3.1).

\subsubsection*{Lema 5.4.}
\underline{Sea $a>0$ un entero y sea
$A=\mathbb Z[[X_1,\ldots,X_a]]$ el anillo de series formales en $a$
indeterminadas. Sea $B$ la $A$-álgebra}
\[
B=A[[T]]\bigg/\left(\prod_{i=1}^{a}(T-X_i)\right).
\]
\underline{Entonces $B$ es un $A$-módulo libre con base
$1,T,\ldots,T^{a-1}$:}
\begin{equation}\tag{5.4.1}
B=AT^{a-1}\oplus AT^{a-2}\oplus\cdots\oplus A.
\end{equation}

\pageoriginale

\underline{Sea $S_i\in A$ la $i$-ésima función simétrica elemental}
\[
S_i=\sum_{1\leq\alpha_1<\cdots<\alpha_i\leq a}
X_{\alpha_1}\cdots X_{\alpha_i},
\]
\underline{de modo que}
\begin{equation}\tag{5.4.2}
\sum_{i=0}^{a}(-1)^iS_iT^{a-i}=0
\qquad\text{en }B.
\end{equation}
\underline{El elemento $1+S_1+\cdots+S_a$ es inversible en $A$;
escribamos su inversa como}
\begin{equation}\tag{5.4.3}
\frac{1}{\sum_{i=0}^{a}S_i}
=\sum_{j\geq0}b_j,
\qquad b_j\text{ homogéneo de grado }j\text{ en }A,
\qquad b_0=1.
\end{equation}
\underline{Para cada entero $j\geq0$, definamos un elemento homogéneo
$a_j\in A$ de grado $j$ escribiendo, mediante (5.4.1),}
\begin{equation}\tag{5.4.4}
T^{a-1+j}
=a_jT^{a-1}+{}?\,T^{a-2}+{}?\,T^{a-3}+\cdots+{}?,
\qquad ?\in A.
\end{equation}
\underline{Entonces, para todo $j\geq0$,}
\begin{equation}\tag{5.4.5}
a_j=(-1)^jb_j.
\end{equation}

\noindent\underline{Demostración.}
En vista de la definición (5.4.3) de los $b_j$, basta demostrar

\pageoriginale

\begin{equation}\tag{5.4.6}
\left\{
\begin{aligned}
a_0&=1,\\
\sum_{i=0}^{a}(-1)^iS_i\,a_{j-i}&=0
&&\text{si }j>0.
\end{aligned}
\right.
\end{equation}
Es evidente que $a_0=1$. Para $j>0$ fijo, la segunda relación se
obtiene tomando el coeficiente de $T^{a-1}$, mediante (5.4.1), en la
relación
\begin{equation}\tag{5.4.7}
\sisi{\displaystyle\sum_{i=0}}{\displaystyle\sum_{i=0}^{a}}
(-1)^iS_iT^{a-1+j-i}=0
\qquad\text{en }B.
\end{equation}
\footnote{La impresión omite la cota superior $a$ en (5.4.7); la
relación es (5.4.2) multiplicada por $T^{j-1}$.}
Esta relación se obtiene multiplicando (5.4.2) por $T^{j-1}$, lo que
demuestra el lema.

\subsubsection*{Corolario 5.5.}
\underline{Se tiene}
\begin{equation}\tag{5.5.1}
\deg(\check X)
=\frac{(-1)^n}{\deg(\varphi)}
\int_X\frac{(1+H)^{r-1}}{C(N)}.
\end{equation}

\noindent\underline{Demostración.}
Se aplica (5.3.3), con $M=N$ el haz normal de
$X\hookrightarrow\mathbb P^r$, a (5.2.8).

Para una variedad cuasiproyectiva lisa $S$, escribamos
\begin{equation}\tag{5.5.2}
C(S)\underset{\mathrm{def}}{=}C(T_S)
\end{equation}
para su clase de Chern total, donde $T_S$ denota el fibrado tangente
de $S$.

\pageoriginale

Dualizando la sucesión exacta (5.1.2), se obtiene
\begin{equation}\tag{5.5.3}
C(\mathbb P^r)=(1+H)^{r+1}
\qquad\text{en }\operatorname{CH}(\mathbb P^r).
\end{equation}
Dualizando la sucesión exacta (5.1.1), se obtiene
\begin{equation}\tag{5.5.4}
C(N)\,C(X)=C(\mathbb P^r)|_X
\qquad\text{en }\operatorname{CH}(X).
\end{equation}
Combinando (5.5.3) y (5.5.4), la fórmula (5.5.1) se convierte en:

\subsubsection*{Corolario 5.6.}
\underline{Se tiene}
\begin{equation}\tag{5.6.1}
\deg(\check X)
=\frac{(-1)^{\sisi{r}{n}}}{\deg(\varphi)}
\int_X\frac{C(X)}{(1+H)^2}.
\end{equation}
\footnote{La impresión lleva $(-1)^r$. La sustitución de (5.5.3) y
(5.5.4) en (5.5.1) no altera su signo $(-1)^n$; la fórmula (5.7.3)
que sigue confirma el mismo exponente.}

\subsubsection*{5.7.}
Para una variedad proyectiva lisa $S$, definimos su
\underline{característica de Euler--Poincaré} por
\begin{equation}\tag{5.7.1}
\chi(S)=\int_SC(S).
\end{equation}

\subsubsection*{Proposición 5.7.2.}
\underline{Sean $Y$ una sección hiperplana lisa de $X$ y $\Delta$ una
sección lisa de $X$ por un subespacio lineal de codimensión dos de
$\mathbb P^r$. Entonces}
\begin{equation}\tag{5.7.3}
\deg(\check X)
=\frac{(-1)^n}{\deg(\varphi)}
\bigl[\chi(X)+\chi(\Delta)-2\chi(Y)\bigr].
\end{equation}

\pageoriginale

\noindent\underline{Demostración.}
Sea $i:Y\hookrightarrow X$ la inclusión. Se tiene
\[
N_{Y/X}=H|_Y=i^*H.
\]
La sucesión exacta
\[
0\longrightarrow T_Y
\longrightarrow i^*T_X
\longrightarrow i^*H
\longrightarrow0
\]
da, por tanto,
\begin{equation}\tag{5.7.4}
C(Y)=i^*\!\left(\frac{C(X)}{1+H}\right).
\end{equation}
De ahí
\[
\chi(Y)=\int_Yi^*\!\left(\frac{C(X)}{1+H}\right),
\]
y la fórmula de proyección da
\begin{equation}\tag{5.7.5}
\chi(Y)
=\int_Xi_*i^*\!\left(\frac{C(X)}{1+H}\right)
=\int_X\frac{H\,C(X)}{1+H}.
\end{equation}
Iterando este cálculo para $\Delta$, se obtiene
\begin{equation}\tag{5.7.6}
\begin{aligned}
\chi(\Delta)
&=\int_Y\frac{H\,C(Y)}{1+H}\\
&=\int_Yi^*\!\left(
  \frac{H}{1+H}\cdot\frac{C(X)}{1+H}
  \right)
=\int_X\frac{H^2C(X)}{(1+H)^2}.
\end{aligned}
\end{equation}
Por consiguiente,
\[
\begin{aligned}
\chi(X)-2\chi(Y)+\chi(\Delta)
&=\int_XC(X)\left[
1-\frac{2H}{1+H}+\frac{H^2}{(1+H)^2}
\right]\\
&=\int_X\frac{C(X)}{(1+H)^2}.
\end{aligned}
\]
Junto con (5.6.1), esto demuestra la proposición.

\pageoriginale

\section*{6. El caso de una base general}

\subsubsection*{6.1.}
Sean $S$ un esquema y $\bar s$ un punto geométrico de $S$, definido
por una clausura algebraica del cuerpo residual $k(s)$ para un punto
$s\in S$. Sean $\mathbb P^r_S$ el espacio proyectivo de dimensión
$r$ sobre $S$, $\check{\mathbb P}^{r}_S$ su espacio proyectivo dual,
y
\[
I\subset\mathbb P^r_S\times_S\check{\mathbb P}^{r}_S
\]
la variedad de incidencia. Sea
$D\subset\check{\mathbb P}^{r}_S$ una recta definida por un punto con
valores en $S$ de la grassmanniana apropiada, y sea $\Delta$ el eje
del pincel de hiperplanos correspondiente (cf. 2.1).

Sea $X$ un $S$-esquema propio y liso, puramente de dimensión
$n\geq1$, provisto de una inmersión proyectiva
$X\hookrightarrow\mathbb P^r_S$. Pongamos
\[
\widetilde X=(X\times_SD)\cap I.
\]
Así, $\widetilde X$ es el esquema sobre $D$ cuyas fibras son las
secciones hiperplanas de $X$ parametrizadas por $D$; sea
$u:\widetilde X\to D$ la proyección.

Utilizamos un subíndice $\bar s$ para denotar el cambio de base por
$\bar s\to S$. Supongamos que $\Delta_{\bar s}$ es transversal a
$X_{\bar s}$. Entonces, sobre un entorno $U$ de $s$ en $S$,
$\Delta$ es transversal a $X$. En lo que sigue reemplazamos $S$ por
este entorno y escribimos $U=S$. El esquema $\widetilde X$ es
entonces liso sobre $S$, y un entorno en $\widetilde X$ de la imagen
inversa de $\Delta$ es liso sobre $D$.

Supongamos además que una de las secciones hiperplanas de
$X_{\bar s}$ pertenecientes al pincel $D_{\bar s}$ es lisa. Tras
reemplazar de nuevo $S$ por un entorno apropiado de $s$, el subesquema
cerrado $F\subset\widetilde X$ definido por el ideal jacobiano
$J^{n-1}(\widetilde X/D)$ es finito y plano sobre $S$.

\pageoriginale

\subsubsection*{6.2.}
El \underline{lugar excepcional} del pincel $D$ es el divisor de
Cartier relativo $E$ de $D/S$ definido por el
$\mathcal O_D$-módulo finito y plano $u_*\mathcal O_F$. Localmente
sobre $D$, si
\[
\mathcal O_D^m\xrightarrow{\ v\ }\mathcal O_D^m
\longrightarrow u_*\mathcal O_F\longrightarrow0
\]
es una presentación de $u_*\mathcal O_F$, entonces $E$ tiene ecuación
$\det(v)=0$. La formación de $E$ conmuta con todo cambio de base.

El divisor $E_{\bar s}$ sobre $D_{\bar s}$ puede escribirse
\[
E_{\bar s}=\sum_Pn(P)[P],
\]
donde $n(P)$ es la suma de las colongitudes del ideal jacobiano
$J^{n-1}(\widetilde X/D)$ en los puntos críticos de $u$ que se
aplican en $P$.

\subsubsection*{Proposición 6.3.}
\underline{Supongamos que $D_{\bar s}$ es un pincel de Lefschetz de
secciones hiperplanas de $X_{\bar s}$ y que
$\operatorname{car}k(s)\ne2$ o que $n$ es par. Existe un entorno $U$
de $s$ en $S$ tal que:}
\begin{enumerate}
\item[(a)] \underline{para todo punto geométrico $\bar t$ de $U$,
$D_{\bar t}$ es un pincel de Lefschetz de secciones hiperplanas de
$X_{\bar t}$;}
\item[(b)] \underline{sobre $U$, el lugar excepcional $E$ es étale
sobre $S$.}
\end{enumerate}
Para demostrar (a) y (b), basta observar que, bajo estas hipótesis,
$D_{\bar s}$ satisface la condición 2.2(c) si y solo si el divisor
$E_{\bar s}$ es reducido. Invocando 3.3, se podría utilizar solamente
la implicación más fácil de izquierda a derecha.

\subsubsection*{6.3 (caso excepcional).}
\footnote{La impresión repite el número 6.3 después de la Proposición
6.3.}
La proposición es falsa en el caso excepcional recién mencionado.
No obstante, si $S$ tiene característica dos y $n$ es impar, existe
un entorno $U$ de $s$ sobre el cual se cumple (a), y se puede definir
un divisor $E'$ étale sobre $U$ tal que $E=2E'$. No utilizaremos este
resultado.

\subsection*{Referencias}

\noindent[1] A. Grothendieck, \emph{La théorie des classes de Chern},
Bull. Soc. Math. France \textbf{86} (1958), 137--154.
